%%%%%%
%
% $Autor: Wings $
% $Datum: 2020-01-18 11:15:45Z $
% $Pfad: WuSt/Skript/Produktspezifikation/powerpoint/ImageProcessing.tex $
% $Version: 4620 $
%
%%%%%%

\chapter{Erzeugung des Trainingsmodells}


In diesem Kapitel werden die wichtigsten Programmabschnitte aus dem ursprünglichen
Programms zur Erzeugung des neuronalen Netzes erläutert. Zu Verfügung gestellt wurde
das Programm von Herrn Dr.-Ing. Ole Mathis Magens.

\section{Trainingsmodell}

Die wichtigste Bibliothek in dem nachfolgenden Programmcode ist die Tensorflow-Bibliothek,
da diese auf die wichtigsten Keras Funktionen zugreift. Der nachfolgende Programmcode
ist an einigen Stellen gekürzt, ein vollständiger Programmcode befindet sich im Anhang~B.

\begin{lstlisting}
import tensorflow as tf
from tensorflow.python.keras.datasets import cifar10
from tensorflow.python.keras.preprocessing.image import ImageDataGenerator
from tensorflow.python.keras.models import Sequential
from tensorflow.python.keras.layers import Dense , Dropout , Activation , Flatten
from tensorflow.python.keras.layers import Conv2D , MaxPooling2D
from tensorflow.python.keras.callbacks import TensorBoard
import collections
\end{lstlisting}

	Die Bibliothek collections ist für die Berechnung der Gewichte zuständig und somit 
	ebenfalls von Relevanz.
	Als erstes werden die Trainingsdaten aus Kapitel 6 als Variable pickle in geöffnet und
	wieder unter dem Variablennamen X, y, Xval, yval gespeichert. Das rb steht für read
	bytes.

\begin{lstlisting}
pickle_in = open("X.pickle" ,"rb")
X = pickle.load( pickle_in )

pickle_in = open("y.pickle" ,"rb")
y = pickle.load( pickle_in )
pickle_in = open("Xval.pickle" ,"rb")
Xval = pickle.load( pickle_in )

pickle_in = open("yval.pickle" ,"rb")
yval = pickle.load( pickle_in )
\end{lstlisting}

Die Byte-Stream Zahlen enthalten Zahlenwerte zwischen 0-255, die sich aus dem entsprechenden 
Farbwert des Pixels ergeben. Das neuronale Netz auf Basis von Tensorflow
arbeitet im Wertebereich zwischen -1 bis 1.[Gmb19a] Die Variablen mit den Bildinformationen 
wird auf einen Wertebereich zwischen 0-1 skaliert.


\begin{lstlisting}
X = X/255.0
Xval = Xval /255.0
\end{lstlisting}

Mit der Funktion ImageDataGeneration können die einzelnen Bilder in ihrer Ausrichtung
und ihrer Helligkeit leicht verändert werden. Hierdurch soll das Netz abstraktere Konzepte
lernen, um auch bei leichten Abweichungen des Bildausschnitts gut zu performen.

\begin{lstlisting}
datagen = ImageDataGenerator (
rotation_range = 5 ,
brightness_range = (0.2 ,1.5) ,
width_shift_range =0.1 ,
height_shift_range =0.1 ,
zoom_range =0.1 ,
shear_range =0.1 ,
fill_mode='nearest')
datagen.fit(X)
\end{lstlisting}

\begin{lstlisting}
model = Sequential ()

model.add( Conv2D (16 , (3 , 3) , input_shape=(IMG_HEIGHT , IMG_WIDTH , 1) , padding='same') )
model.add( Activation ('relu') )
model.add( MaxPooling2D ( pool_size=(2, 2) ) )

model.add( Conv2D (32 , (3 , 3) , padding='same') )
model.add( Activation ('relu') )
model.add( MaxPooling2D ( pool_size=(2, 2) ) )

model.add( Conv2D (32 , (3 , 3) , padding='same') )
model.add( Activation ('relu') )
model.add( MaxPooling2D ( pool_size=(2, 2) ) )

model.add( Flatten () )
model.add( Dense (32) )
model.add( Activation ('relu') )
model.add( Dropout (0.5) )

model.add( Dense (1) )
model.add( Activation ('sigmoid') )
model.compile (loss='binary_crossentropy' ,
optimizer='AdaDelta' , metrics=['accuracy' ])

model.fit_generator ( datagen.flow(X , y , batch_size=64) ,
epochs=40,
steps_per_epoch=200,
validation_data=datagen.flow(Xval , yval , batch_size=64))
\end{lstlisting}

Der wichtigste Teil des Netzes liegt in dem Aufbau der Netzstruktur. Zuerst wird ein
Modell vom Typ Sequential aufgerufen. Bei dieser Art handelt es sich um einen linea-
ren Stapel von Schichten. Es folgen eine Reihe von Ebenenfunktionen. Ein \ac{cnn} beginnt
mit einer Faltungsschicht (siehe Kapitel 4), welches mit dem Befehl model.add(Conv2D)
erzeugt wird. Es folgt in Klammern die Anzahl der Filter (hier 16) und die Größe des
Filters (hier: 3x3), danach wird die Eingangsgröße der Trainingsdaten und die Art des
padding mitgeteilt. Mit der Funktion model.add(Activation) kann die Aktivierungsfunktion 
der entsprechenden Schicht festgelegt werden. Danach folgt eine Polling-Schicht, 
welche mit dem Befehl model.add(MaxPooling2D)erzeugt wird. In den Klammer befindet
sich die Größe des Kernel (hier: 2x2). Diese Schichtenpakete wiederholen sich zweimal.
Bei der zweiten Faltungsschicht wurde die Filteranzahl auf 32 erhöht. Nach diesem Aufbau 
werden die 2-dimensionalen Bildinhalte mit dem Befehl model.add(Flatten) in einen
eindimensionalen Vektor konvertiert. Anschließend wird ein Dense-Layer mit dem Befehl
model.add(Dense) hinzugefügt. Diese letzte Schicht des \ac{cnn} schafft den Übergang zur
klassischen neuronalen Netzstruktur, in der alle Neuronen untereinander vollverbunden
sind. An dieser Stelle kann ein Dropout von 50\% hinzugefügt werden. In der letzten
Schicht des klassischen neuronalen Netzes gibt es erneut einen Dense-Layer mit der 
Aktivierungsfunktion Sigmoid. Bei binären Klassen wird der Dense-Layer auf 1 gesetzt.
Im nächsten Schritt wird der Compiler definiert. Binary crossentropy wird verwendet,
wenn nur zwei Beschriftungsklassen (0 und 1) vorhanden sind. Mit dem Optimierer Ada-Delta 
wird eine stochastische Gradientenabstiegsmethode verwendet, um den kontinuierlichen Rückgang 
der Lernrate während des Trainings und die Notwendigkeit einer manuell
ausgewählten globalen Lernrate zu umgehen. [Cho19] Die Metrik accuracy berechnet, wie
oft Vorhersagen mit der Beschriftung aus den Kategorien übereinstimmen. Mit 
dem model.fit Generator werden wichtige Einstellungen zum Training definiert. Es werden die
Trainings und Validierungsdaten als daten.flow zugeordnet, sowie dem Modell die Epo-
chenanzahl, Stapelgröße und Anzahl der Schritte. In der Variable NAME kann ein eigens
definierter, geeigneter Name für das neuronale Netz angegeben werden. Es empfiehlt sich,
dem Modell Informationen der Netzstruktur mitzugeben.

\begin{lstlisting}
NAME = "demold_16_32_32 -32-32CNN-aug -50drop-ver180 -single_mold"
model.save(NAME+'.model')
\end{lstlisting}

\section{Vorhersage}

Nachdem das neuronale Netz trainiert wurde, geht es in diesem Programmabschnitt darum 
die Performance des neuronalen Netzes in der Realität zu überprüfen. Das Modell
wird dabei aktiv auf Bilder aus bestimmten Ordnern angewendet, um zu testen welcher
Kategorie die Bilder zugeordnet werden.

\begin{lstlisting}
CATEGORIES = [ "good" , "bad" ]

files = glob.glob('./Testmenge_good_30x80/*.bmp')

def prepare ( filepath ) :
    IMG_HEIGHT = 80
    IMG_WIDTH = 30
    img_array = cv2.imread (filepath , cv2.IMREAD_GRAYSCALE )
    new_array = cv2.resize (img_array , (IMG_WIDTH , IMG_HEIGHT ) )
        return new_array.reshape ( - 1 , IMG_HEIGHT , IMG_WIDTH ,1)

model = tf.keras.models.load_model ('Model_300ValSam_30x80_E40_S200_FE3_DrO50Pr.model')

prediction = numpy.zeros (( len( files ) ,1 , 1) )
\end{lstlisting}

Das Programm nutzt die Bibliotheken cv2, os, tensorflow, matplotlib, glob und numpy.
Nach dem Einspielen der Bibliotheken werden die möglichen Kategorien definiert. Der
Index 0 definiert die Katergorie good und der Index 1 die Kategorie bad. Die Variable files
enthält den Ordnerpfad, in dem sich die Bilder befinden, die getestet werden sollen. Der
Pfad aus dem Programmabschnitt führt zu einem Ordner mit Testbildern der Kategorie
good.

Die Funktion prepare bereitet die Bilder aus dem Ordner so vor, dass sie sich in einem
Format befinden, welches das neuronale Netz verwenden kann. Diese Funktion ist ähnlich
der Funktion zur Erstellung der Trainingsdaten aus Kapitel 6. Anschließend wird das
neuronale Netz in der Variable model als neues Array über die angegeben Ordnerlänge
gespeichert. Die for-Schleife erzeugt eine Laufvariable i, die über die gesamte Ordnerlänge
des Testbilder aus der Kategorie good läuft. Die Funktion model.predict übernimmt die
Klassifikationen des neuronalen Netzes und prüft die Bilder in dem Testordner auf deren 
Klasse. In der if-Schleife werden alle Bilder angezeigt, die nicht der Kategorie good
zugeordnet werden können. Da sich in dem Testordner nur Bilder der Kategorie good
befinden, weisen alle angezeigten Bilder die der Kategorie bad zugeordnet wurden, auf
eine Ungenauigkeit des Netzes hin.

\begin{lstlisting}
for i in range (len( files ) ) :
    prediction [ i ,0 ,0] = model.predict ([ prepare ( files [ i ]) ])
    print ( files [ i]+': '+CATEGORIES [ int( prediction [ i ] [ 0 ] [ 0 ] ) ])
    img_array = cv2.imread ( files [ i ] ,cv2.IMREAD_COLOR )
    pyplot.imshow (img_array , cmap='gray')
    pyplot.show ()

if CATEGORIES [ int( prediction [ i ] [ 0 ] [ 0 ] ) ] == 'bad' :
    pyplot.imshow (img_array , cmap='gray')
    pyplot.show ()
\end{lstlisting}

